% Descent of quasi-coherent sheaves along universal effective epimorphisms of schemes — standalone article for Hexagon
% (question of the Buffalo lecture of 11 July 1973, video 3ULFIl_2Wis, 6:26:20–6:27:24), from Appendix C of docs/buffalo-open.
% pdfLaTeX. Compile: tectonic -X compile article.tex.
\documentclass[11pt,a4paper]{amsart}
\usepackage[T1]{fontenc}
\usepackage[utf8]{inputenc}
\usepackage{lmodern,microtype,amssymb,mathtools}
\usepackage[margin=2.7cm]{geometry}
\usepackage[hidelinks]{hyperref}

\newtheorem{theorem}{Theorem}[section]
\newtheorem{proposition}[theorem]{Proposition}
\newtheorem{lemma}[theorem]{Lemma}
\newtheorem{corollary}[theorem]{Corollary}
\newtheorem*{theoremA}{Theorem A}
\newtheorem*{theoremB}{Theorem B}
\newtheorem*{theoremC}{Theorem C}
\theoremstyle{definition}
\newtheorem{definition}[theorem]{Definition}
\theoremstyle{remark}
\newtheorem{remark}[theorem]{Remark}

\newcommand{\q}[1]{\guillemotleft\,#1\,\guillemotright}
\newcommand{\Z}{\mathbf Z}
\newcommand{\Q}{\mathbf Q}
\newcommand{\Pp}{\mathcal P}
\newcommand{\QCoh}{\mathrm{QCoh}}
\newcommand{\Des}{\mathrm{Des}}
\newcommand{\Spec}{\operatorname{Spec}}
\newcommand{\Hom}{\operatorname{Hom}}
\newcommand{\Mor}{\operatorname{Mor}}

\title[Descent along universal effective epimorphisms]{Descent of quasi-coherent sheaves\\ along universal effective epimorphisms of schemes}
\author{Michel Hua}
\address{Commutator, Paris, France}
\email{michel@commutator.io}
\thanks{ORCID 0009-0002-8304-0683.}
\date{8 November 2026}
\subjclass[2020]{14A15; 18F10, 14F06, 13C11, 54D40, 03E50}
\keywords{descent, quasi-coherent sheaves, canonical topology, universal effective epimorphisms, pure morphisms, Stone--\v Cech remainder, butterfly points, Grothendieck}

\begin{document}
\begin{abstract}
In a lecture given at Buffalo in 1973, Grothendieck asked whether quasi-coherent sheaves descend along universal effective epimorphisms of schemes, the single-arrow coverings of the canonical topology, and whether pulling them back is at least fully faithful; he doubted it could be proved. We show that pulling back is fully faithful along every universal effective epimorphism, with no quasi-compactness assumption. Along a quasi-compact one, descent is effective; this is due to Mesablishvili, and we give a short independent proof that reduces to the affine case, where universal effective epimorphisms are the pure ring maps. Without quasi-compactness effectivity can fail. Suppose that some nonprincipal ultrafilter on the primes is generated by $\aleph_1$ sets, as every one is under the continuum hypothesis. Then the universal effective epimorphism that the Stacks Project uses to show that $\Spec\Z$ is not quasi-compact in the canonical topology carries a descent datum, an extension of the structure sheaf by $\bigoplus_p\Z/p$, that is not effective. So, under that hypothesis, quasi-coherent sheaves form a separated prestack that is not a stack for the canonical topology, on schemes and already on affine schemes. Along this covering, effectivity is equivalent to the triviality of certain coherent families, a condition on the point of $\beta\mathbf N\setminus\mathbf N$ defined by the ultrafilter. For $\{0,1\}$-valued families it says that the point is not a strong butterfly point in the sense of Fine and Gillman. In the model of van Douwen, Kunen and van Mill no point is one, and the question in ZFC remains open. The affine core of full faithfulness is checked in Lean.
\end{abstract}
\maketitle

\section{The question}

In a category with fibre products, a morphism $f\colon X\to Y$ is an \emph{effective epimorphism} if, for every object $T$,
\[
\Mor(Y,T)\to\Mor(X,T)\rightrightarrows\Mor(X\times_YX,T)
\]
is an equaliser, and a \emph{universal effective epimorphism} if every base change of $f$ is one. The \emph{canonical topology} is the finest Grothendieck topology for which every representable functor is a sheaf. Its covering families are the universally effective-epimorphic families, and a sieve generated by a single arrow covers if and only if that arrow is a universal effective epimorphism \cite[00WP, 00ZA]{stacks}, \cite[\S1.2]{af}. In a category with fibre products, effective, regular and strict epimorphisms coincide; Mesablishvili calls universal effective epimorphisms \emph{stable regular epimorphisms}. In the category of schemes, coproducts are disjoint and universal, so a family $\{X_i\to Y\}$ is universally effective-epimorphic if and only if the single morphism $\coprod_iX_i\to Y$ is a universal effective epimorphism. Every fpqc covering is one \cite[023Q]{stacks}.

For $f\colon X\to Y$, write $\Des(f)$ for the category of quasi-coherent $\mathcal O_X$-modules with descent data relative to $f$, and $f^*\colon\QCoh(Y)\to\Des(f)$ for the pull-back. In the vocabulary of Grothendieck's lectures, $f$ is a \emph{descent morphism} for quasi-coherent sheaves if $f^*$ is fully faithful, and an \emph{effective descent morphism} if it is an equivalence. Faithfully flat quasi-compact morphisms are effective descent morphisms \cite[023T]{stacks}; without quasi-compactness even full faithfulness can fail for faithfully flat morphisms, for instance for the disjoint union of the spectra of the local rings of the complex affine line at its closed points, an example given on MathOverflow \cite{mo127362}.

In his algebraic geometry course at the State University of New York at Buffalo, at the session of 11 July 1973, Grothendieck proved that schemes are sheaves for the fpqc topology. A listener asked why one does not take as coverings all the universal effective epimorphisms. He answered that this gives the finest topology for which representable functors are sheaves, but that such morphisms are hard to recognise, and that one does not know whether quasi-coherent sheaves descend along them, or even whether pulling them back is faithful; he added that he doubted \q{very much} that this could be proved \cite[6:26:20--6:27:24]{buffalo}. The recording is transcribed from captions and by machine; the transcription has not been checked by ear, and only these two words are quoted.

Both questions are answered below, and the answer on full faithfulness goes against his doubt.

\begin{theoremA}[Theorem~\ref{t:A}]
Every universal effective epimorphism $f\colon X\to Y$ of schemes is a descent morphism for quasi-coherent sheaves: $f^*\colon\QCoh(Y)\to\Des(f)$ is fully faithful. No quasi-compactness is needed.
\end{theoremA}

\begin{theoremB}[Mesablishvili; Theorem~\ref{t:B}]
Every quasi-compact universal effective epimorphism of schemes is an effective descent morphism for quasi-coherent sheaves.
\end{theoremB}

\begin{theoremC}[Theorem~\ref{t:C}]
Suppose that some nonprincipal ultrafilter $U$ on the set $\Pp$ of primes is generated by $\aleph_1$ sets, as is every one under the continuum hypothesis. For $A\in U$ let $\Z_A$ be $\Z$ with the primes of $A$ inverted. Then $f\colon\coprod_{A\in U}\Spec\Z_A\to\Spec\Z$ is a universal effective epimorphism of schemes, hence a descent morphism, but it is not an effective descent morphism for quasi-coherent sheaves. Consequently $\QCoh$ is a separated prestack, but not a stack, for the canonical topology on schemes, and also for the canonical topology on affine schemes.
\end{theoremC}

The covering of Theorem~C is the one with which the Stacks Project shows that $\Spec\Z$ is not quasi-compact in the canonical topology \cite[0EUE, 0EUG]{stacks}. Its descent data are analysed in Section~\ref{sec:nqc}: effectivity along it is equivalent to the triviality of certain coherent families of sections (Proposition~\ref{p:dev}), and for $\{0,1\}$-valued families this triviality says that the point $U$ of $\omega^*=\beta\Pp\setminus\Pp$ is not a strong butterfly point, in the sense of Fine and Gillman \cite{fg}. At a point of character $\aleph_1$ the butterfly lemma of Fine and Gillman, in the form given by Pitz and Suabedissen \cite[Lemma~3.1]{pitz}, produces a strong butterfly point, hence a non-effective datum. In the model of van Douwen, Kunen and van Mill \cite{dkm} there are no strong butterfly points in $\omega^*$, and we do not know what happens there (Section~\ref{sec:open}).

\subsection*{What is in print}
We searched zbMATH Open, arXiv, Crossref, Semantic Scholar, MathOverflow and the Stacks Project, and read the following. The affine case is classical: for a ring map $A\to B$, being a universal effective epimorphism, being pure, being a descent morphism and being an effective descent morphism for modules are equivalent. This is due to Joyal and Tierney (unpublished), Olivier \cite{olivier}, Mesablishvili \cite{mesablishvili}, and is in the Stacks Project \cite[08WE, 08XA, 0EUA]{stacks}; André and Fiorot give a full account, with a reconstruction of Olivier's proof \cite[\S\S2--3]{af}. Mesablishvili proved, for quasi-compact morphisms of schemes, that pure, stable regular epimorphism and stable effective descent morphism for quasi-coherent sheaves are equivalent \cite{mesablishvili-acs,mesablishvili2004}; this contains Theorem~B, and Theorem~A in the quasi-compact case. Rydh's study of submersive and subtrusive morphisms \cite{rydh}, Picavet's \cite{picavet}, and Kern's recent crisp topology \cite{kern}, which extends purity beyond quasi-compact morphisms, do not treat the descent of quasi-coherent sheaves along non-quasi-compact universal effective epimorphisms. Orlov studies quasi-coherent sheaves on presheaves on affine schemes and the effective descent topology \cite{orlov}. Aoki and Burklund observe that the canonical topology on affine schemes is not finitary \cite[footnote~1]{ab}; the covering of Theorem~C is not finitary either. As far as this search went, Theorem~A without quasi-compactness, Theorem~C and the analysis of Section~\ref{sec:nqc} are not in print; we do not claim more than that.

Theorem~C bears on one sentence of André and Fiorot. They prove that the canonical topology on affine schemes coincides with the effective descent topology \cite[Thm.~3.2]{af}, treating a covering given by a single ring map, and note that the argument can be adapted to a general family by taking products. For a finite family this is clear, since the family is equivalent to the single map to the product. Under the hypothesis of Theorem~C the adaptation fails for infinite families. By \cite[0EUA]{stacks} and the equaliser condition of \cite[0EUF]{stacks}, the family $\{\Spec\Z_A\to\Spec\Z\}_{A\in U}$ covers $\Spec\Z$ in the canonical topology of affine schemes, but it is not an effective descent family; so on affine schemes over $\Z$ the two topologies differ. The two coincide on the coverings given by single morphisms, and by finite families.

\section{The affine case}\label{sec:affine}

\begin{proposition}\label{p:affine}
For a ring map $A\to B$ the following are equivalent:
\begin{enumerate}
\item[(a)] $\Spec B\to\Spec A$ is a universal effective epimorphism of schemes;
\item[(b)] it is a universal effective epimorphism of affine schemes;
\item[(c)] for every $A$-module $M$, the sequence $M\to M\otimes_AB\rightrightarrows M\otimes_AB\otimes_AB$ is an equaliser;
\item[(d)] $A\to B$ is pure, that is, universally injective;
\item[(e)] $A\to B$ is an effective descent morphism for modules;
\item[(f)] $A\to B$ is a descent morphism for modules.
\end{enumerate}
\end{proposition}

\begin{proof}
(a)$\Rightarrow$(b): affine schemes form a full subcategory of schemes closed under fibre products, so base changes and test objects can be restricted to affine ones. (b)$\Leftrightarrow$(c) is \cite[0EUA]{stacks}, applied to a family with one member. Its proof of (b)$\Rightarrow$(c) uses the base change along $\Spec(A\oplus M)$, the square-zero extension, and the test object $\mathbf A^1$; the proof of Theorem~\ref{t:A} below runs the same argument over a non-affine base. (c)$\Rightarrow$(d) is clear, as the first map of (c) is injective for every $M$. (d)$\Leftrightarrow$(e)$\Leftrightarrow$(f) is \cite[08XA]{stacks}, a theorem of Joyal and Tierney, announced by Olivier \cite{olivier}, of which Mesablishvili gives a published proof \cite{mesablishvili} and André and Fiorot another \cite[Prop.~3.3]{af}. (d)$\Rightarrow$(a): a pure morphism of affine schemes is a regular epimorphism in the category of schemes \cite[Thm.~5.7, Cor.~5.8]{mesablishvili2004}, and this is stable under base change \cite[Cor.~5.10]{mesablishvili2004}.
\end{proof}

So the affine universal effective epimorphisms are exactly the pure ring maps, whether one tests against affine schemes or against all schemes. Mesablishvili proves more: a quasi-compact morphism of schemes is pure, in the sense that $\mathcal O_Y\to f_*\mathcal O_X$ stays injective after every base change, if and only if it is a stable regular epimorphism \cite[Thm.~5.11]{mesablishvili2004}, and if and only if it is a stable effective descent morphism for quasi-coherent sheaves \cite{mesablishvili-acs}.

\section{Full faithfulness}\label{sec:A}

\begin{theorem}\label{t:A}
Every universal effective epimorphism $f\colon X\to Y$ of schemes is a descent morphism for quasi-coherent sheaves.
\end{theorem}

\begin{proof}
\emph{Step 1: reduction to an affine base.} The question is Zariski-local on $Y$. For quasi-coherent $\mathcal O_Y$-modules $M$ and $F$, morphisms $M\to F$ form a sheaf on $Y$, and so do the morphisms $f^*M\to f^*F$ of descent data, restricted over the open subschemes of $Y$. The base change of $f$ to an open subscheme of $Y$ is again a universal effective epimorphism. If $f^*$ is fully faithful over every affine open subscheme of $Y$, the two sheaves of morphisms agree on a basis of the topology of $Y$, hence on $Y$. So we assume $Y=\Spec A$.

\emph{Step 2: the square-zero extension.} Let $F$ be a quasi-coherent $\mathcal O_Y$-module, and let $Y'=\Spec_Y(\mathcal O_Y\oplus F)$, the square-zero extension of $Y$ by $F$; it is affine, with ring $A\oplus\Gamma(Y,F)$. Put $X'=X\times_YY'=\Spec_X(\mathcal O_X\oplus f^*F)$. Its underlying space is $|X|$, and as sheaves on $|X|$ we have $\mathcal O_{X'}=\mathcal O_X\oplus f^*F$. Hence
\[
\Gamma(X',\mathcal O_{X'})=\Gamma(X,\mathcal O_X)\oplus\Gamma(X,f^*F),
\]
and likewise $X'\times_{Y'}X'=(X\times_YX)\times_YY'$ has ring of global sections $\Gamma(X\times_YX,\mathcal O)\oplus\Gamma(X\times_YX,p^*F)$, where $p\colon X\times_YX\to Y$ is the structure map. This needs no affineness of $X$.

\emph{Step 3: the test object $\mathbf A^1$.} The morphism $X'\to Y'$ is a base change of $f$, hence an effective epimorphism. Test it against $\mathbf A^1=\Spec\Z[t]$, for which $\Mor(T,\mathbf A^1)=\Gamma(T,\mathcal O_T)$ for every scheme $T$. The two maps $\Gamma(X',\mathcal O)\rightrightarrows\Gamma(X'\times_{Y'}X',\mathcal O)$ are induced by the two projections, which are compatible with the decompositions of Step~2, since the projections are morphisms over $Y'$ and the decompositions come from the $\mathcal O_Y$-algebra $\mathcal O_Y\oplus F$. Taking the components in $F$ gives an equaliser
\[
\Gamma(Y,F)\to\Gamma(X,f^*F)\rightrightarrows\Gamma(X\times_YX,p^*F).
\]

\emph{Step 4: the case $M=\mathcal O_Y$.} A morphism $\mathcal O_X\to f^*F$ is a section $s\in\Gamma(X,f^*F)$, and it is compatible with the canonical descent data exactly when its two pull-backs to $X\times_YX$ agree. So the equaliser of Step~3 says that
\[
\Hom_Y(\mathcal O_Y,F)\to\Hom_{\Des(f)}(f^*\mathcal O_Y,f^*F)
\]
is bijective.

\emph{Step 5: every $M$.} The category $\Des(f)$ has colimits, computed on $X$ with the induced descent data, since pulling back along the projections preserves colimits; and $f^*$ preserves colimits. So both $\Hom_Y(-,F)$ and $\Hom_{\Des(f)}(f^*-,f^*F)$ turn direct sums and cokernels of $\mathcal O_Y$-modules into products and kernels. Since $Y$ is affine, every quasi-coherent $M$ has a presentation $\mathcal O_Y^{(I)}\to\mathcal O_Y^{(J)}\to M\to0$, and Step~4 gives the bijection for $M$.
\end{proof}

In particular $\QCoh$ is a separated prestack for the canonical topology: for every covering family the comparison functor to descent data is fully faithful, the family being replaced by the coproduct of its members.

\subsection{The affine core in Lean}\label{sec:lean}

The affine case of Theorem~\ref{t:A} is checked in Lean~4 with Mathlib, in the file \texttt{lean/\allowbreak Buffalo/\allowbreak QcohDescent.lean} of the repository \cite{repo}, added in commit \texttt{379ce37e}; the proofs use the standard axioms only. Let $A\to B$ be a map of commutative rings and $M$, $N$ modules over $A$. A morphism of descent data $B\otimes_AM\to B\otimes_AN$ is a $B$-linear map $g$ whose two base changes to $B\otimes_AB$, along $b\mapsto b\otimes1$ and $b\mapsto1\otimes b$, agree. The file proves:
\begin{itemize}
\item \texttt{descend\_bijective}: if $N\to B\otimes_AN\rightrightarrows B\otimes_AB\otimes_AN$ is an equaliser, base change
\[\Hom_A(M,N)\to\Hom_{\Des}(B\otimes_AM,\,B\otimes_AN)\]
is bijective for every $A$-module $M$;
\item \texttt{isEqualizer\_iff\_bijective}: conversely, bijectivity for $M=A$ alone gives the equaliser condition for $N$;
\item \texttt{equalizerCondition\_iff\_isPure}: the equaliser condition for all modules is equivalent to $A\to B$ being pure \cite[08WE]{stacks}, and \texttt{descend\_bijective\_of\_isPure} states full faithfulness for a pure map.
\end{itemize}
The formal proof needs no presentation of $M$. A $B$-linear map $g$ is determined by the $A$-linear map $m\mapsto g(1\otimes m)$, and it is compatible with the descent data exactly when this map lands in the equaliser of $B\otimes_AN\rightrightarrows B\otimes_AB\otimes_AN$, which is $N$. Not formalised are: the step from universal effective epimorphisms to the equaliser condition \cite[0EUA]{stacks}, since schemes, effective epimorphisms and the test object $\mathbf A^1$ are not modelled; the effectivity of descent along a pure map \cite[08XA]{stacks}; the non-affine case of Theorem~\ref{t:A}; and Theorems~\ref{t:B} and~\ref{t:C}.

\section{The quasi-compact case}\label{sec:B}

\begin{theorem}[Mesablishvili \cite{mesablishvili-acs,mesablishvili2004}]\label{t:B}
Every quasi-compact universal effective epimorphism $f\colon X\to Y$ of schemes is an effective descent morphism for quasi-coherent sheaves.
\end{theorem}

The theorem follows from Mesablishvili's two characterisations of quasi-compact pure morphisms recalled at the end of Section~\ref{sec:affine}. The following proof is independent of his; it reduces to the affine case by a finite affine cover.

\begin{proof}
\emph{Step 1: reduction to an affine base.} By Theorem~\ref{t:A}, $f^*$ is fully faithful over every open subscheme of $Y$. Effectivity is then Zariski-local on $Y$: if the restrictions of a descent datum over the members of an open cover of $Y$ descend, the descended modules are unique up to unique isomorphism, so they glue. This is the argument Grothendieck gives in the course for fpqc morphisms. So assume $Y=\Spec A$.

\emph{Step 2: a finite affine cover.} As $f$ is quasi-compact, $X$ is quasi-compact. Take a finite affine open cover $(U_i)$ of $X$, put $U=\coprod_iU_i=\Spec B$, and let $q\colon U\to X$ be the canonical map.

\emph{Step 3: the composite is pure.} The map $q$ is a Zariski covering, hence a universal effective epimorphism. Let $R$ be the sieve on $Y$ generated by $g=f\circ q$ and $S$ the sieve generated by $f$; $S$ covers in the canonical topology. For an arrow $h=f\circ k$ of $S$, with $k\colon T\to X$, the pull-back of $R$ along $h$ contains the pull-back along $k$ of the sieve generated by $q$, which covers $T$. By the transitivity axiom of a topology, $R$ covers. So $g$ is a universal effective epimorphism, and $A\to B$ is pure by Proposition~\ref{p:affine}.

\emph{Step 4: descent along $g$.} Let $(N,\theta)$ be a descent datum on $X$ relative to $f$. Its pull-back $q^*(N,\theta)$ is a descent datum relative to $g$, which descends by \cite[08XA]{stacks}: there are an $A$-module $M$ and an isomorphism $\psi\colon g^*M\cong q^*N$ compatible with the descent data on $U\times_YU$.

\emph{Step 5: descent of $\psi$ along $q$.} The morphism $U\times_XU\to U\times_YU$ is an immersion, the base change of the diagonal $X\to X\times_YX$; restrict the compatibility of $\psi$ along it. Along this diagonal, $\theta$ restricts to the identity, as one sees by pulling the cocycle condition back along the diagonal of $X\times_YX\times_YX$. Hence the two pull-backs of $\psi$ to $U\times_XU$ agree: $\psi$ is compatible with the canonical descent data for the Zariski covering $q$, and it descends to an isomorphism $\psi_0\colon f^*M\cong N$ on $X$.

\emph{Step 6: compatibility of $\psi_0$ with $\theta$.} Both $\psi_0$ and $\theta$ live on $X\times_YX$, and it suffices to compare them after pulling back along $U\times_YU\to X\times_YX$. This map is a Zariski covering, by the opens $U_i\times_YU_j$, so pulling back along it is faithful; and after the pull-back, the compatibility is that of $\psi$ in Step~4. So $(N,\theta)\cong f^*M$.
\end{proof}

\begin{remark}\label{r:qcopen}
The proof gives more: if, locally on $Y$, some quasi-compact open $X_0\subseteq X$ is such that $X_0\to Y$ is a universal effective epimorphism, then $f$ is an effective descent morphism. To see it, assume $Y$ affine. A datum $(N,\theta)$ on $X$ restricts to one on $X_0$, which descends by Theorem~\ref{t:B} to some $M$, with $\psi_0\colon f^*M|_{X_0}\cong N|_{X_0}$. Let $\mathrm{pr}_1\colon X\times_YX_0\to X$ be the first projection, a base change of $X_0\to Y$, hence a universal effective epimorphism. On $X\times_YX_0$, $\theta$ followed by the pull-back of $\psi_0^{-1}$ gives an isomorphism $\mathrm{pr}_1^*N\cong\mathrm{pr}_1^*f^*M$, which the cocycle condition for $\theta$ and the compatibility of $\psi_0$ make a morphism of descent data relative to $\mathrm{pr}_1$. By Theorem~\ref{t:A} for $\mathrm{pr}_1$, it and its inverse descend, to an isomorphism $N\cong f^*M$, whose compatibility with the descent data is checked after the faithful pull-back to $X\times_YX\times_YX_0$. This is the statement of the course that a morphism through which an effective descent morphism factors is itself one.
\end{remark}

The next lemma is not used later. It explains why a counterexample must be built from modules that are not Mittag-Leffler, such as localisations, which is the shape of the example of Section~\ref{sec:nqc}.

\begin{lemma}\label{l:ml}
Let $(A\to B_j)_{j\in J}$ be a family of ring maps such that $N\to\prod_jN\otimes_AB_j$ is injective for every $A$-module $N$. If every $B_j$ is a Mittag-Leffler $A$-module, for instance finitely presented, then some finite subfamily is pure: $N\to\prod_{j\in F}N\otimes_AB_j$ is injective for every $N$, for some finite $F\subseteq J$.
\end{lemma}

\begin{proof}
Otherwise, choose for every finite $F\subseteq J$ a module $N_F$ and an element $n_F\neq0$ of $N_F$ that is killed in $N_F\otimes_AB_j$ for every $j\in F$. Choose an ultrafilter $\mathcal U$ on the set of finite subsets of $J$ containing each set $\{F:j\in F\}$, and let $N=\prod_{\mathcal U}N_F$ be the ultraproduct, with $n=[n_F]\neq0$. Fix $j$ and let $Z=\{F:j\in F\}\in\mathcal U$. The ultraproduct is the filtered colimit of the products $\prod_{F\in Z'}N_F$ over $Z'\in\mathcal U$, and tensor products commute with filtered colimits. As $B_j$ is Mittag-Leffler, $(\prod_{F\in Z}N_F)\otimes_AB_j\to\prod_{F\in Z}(N_F\otimes_AB_j)$ is injective; the image of $(n_F)_{F\in Z}\otimes1$ is $(n_F\otimes1)_{F\in Z}=0$, so $(n_F)_{F\in Z}\otimes1=0$, and $n\otimes1=0$ in $N\otimes_AB_j$. So $n$ dies in every $N\otimes_AB_j$, although $n\neq0$, a contradiction with the hypothesis.
\end{proof}

\section{A non-quasi-compact covering}\label{sec:nqc}

\subsection{The covering}
The following is in the Stacks Project \cite[0EUE]{stacks}. Let $\Pp$ be the set of primes and $U$ a nonprincipal ultrafilter on $\Pp$. For $A\in U$ let $S_A$ be the multiplicative set generated by $A$ and $\Z_A=S_A^{-1}\Z$, so that $\Spec\Z_A=\{(0)\}\cup A^c$, where $A^c=\Pp\setminus A$. Then $\Z_A\otimes\Z_B=\Z_{A\cup B}$, $\Z_{\{p\}^c}=\Z_{(p)}$ and $\Z_\Pp=\Q$; for a $\Z$-module $M$ write $M_A=M\otimes\Z_A$. For every $\Z$-module $M$, $M$ is the equaliser of $\prod_{A\in U}M_A\rightrightarrows\prod_{A,B\in U}M_{A\cup B}$ \cite[0EUF]{stacks}, and closed subsets of the $\Spec\Z_A$ that agree on overlaps glue. With \cite[0EUD]{stacks}, by which this equaliser condition together with universal submersiveness gives a universal effective epimorphism of schemes (on submersive morphisms see \cite[\S1]{rydh}), it follows \cite[0EUG]{stacks} that
\[
\mathcal W_U=\{\Spec\Z_A\to\Spec\Z\}_{A\in U},\quad\text{equivalently}\quad f\colon X=\coprod_{A\in U}\Spec\Z_A\to Y=\Spec\Z,
\]
is a universal effective epimorphism of schemes. No finite subfamily is one, since every finite subfamily factors through some $\Spec\Z[1/p]$. So $f$ is flat, surjective, not quasi-compact, and a covering for the canonical topology. Remark~\ref{r:qcopen} does not apply: every quasi-compact open of $X$ meets finitely many components.

\subsection{Descent data}
Since $X\times_YX=\coprod_{A,B}\Spec\Z_{A\cup B}$ and $\QCoh$ of a coproduct is the product of the categories of modules, a descent datum along $f$ consists of $\Z_A$-modules $N_A$ for $A\in U$, isomorphisms $\varphi_{AB}\colon(N_A)_{A\cup B}\cong(N_B)_{A\cup B}$, and the cocycle condition $\varphi_{AC}=\varphi_{BC}\varphi_{AB}$ over $\Z_{A\cup B\cup C}$. This is also the category of descent data for the family $\mathcal W_U$ in the category of affine schemes.

\begin{lemma}\label{l:split}
Let $U$ be an ultrafilter on a set $\alpha$ and $T\subseteq\alpha$. If $T\cap A^c$ is finite for every $A\in U$, then $T$ is finite.
\end{lemma}

\begin{proof}
If $T$ is infinite, split it into two infinite halves $T=T_1\sqcup T_2$. If $T_1^c\in U$, then $T\cap T_1=T_1$ is infinite. Otherwise $T_1\in U$, so $T_2^c\supseteq T_1$ lies in $U$, and $T\cap T_2=T_2$ is infinite.
\end{proof}

\begin{proposition}\label{p:tt}
In ZFC, descent data along $f$ with all $N_A$ torsion, or with all $N_A$ torsion-free, are effective.
\end{proposition}

\begin{proof}
\emph{Torsion data.} If every $N_A$ is torsion, then $N_A=\bigoplus_{p\notin A}T_p$ with $T_p=N_{\{p\}^c}[p^\infty]$, the $p$-primary parts being identified through the $\varphi$'s, and $M=\bigoplus_pT_p$ descends the datum.

\emph{Torsion-free data.} Put $V=N_\Pp$, a $\Q$-vector space ($\Pp\in U$), and identify each $N_A$ with a submodule of $V$ through the $\varphi$'s; write $N_{(p)}=N_{\{p\}^c}\subseteq V$, which makes sense as $\{p\}^c\in U$, $U$ being nonprincipal. Over the principal ideal domain $\Z_A$, the torsion-free module $N_A$ is $\bigcap_{p\notin A}N_{(p)}$. Put $M=\bigcap_pN_{(p)}$. Then $M\subseteq N_A$, so $M_A\subseteq N_A$. Conversely, fix $x\in V$ and let $T(x)=\{p:x\notin N_{(p)}\}$. For every $B\in U$, write $x=y/s$ with $y\in N_B$ and $s\in\Z\setminus\{0\}$; then $T(x)\cap B^c\subseteq\{p:p\mid s\}$, which is finite. By Lemma~\ref{l:split}, $T(x)$ is finite. For $x\in N_A$ we have $T(x)\subseteq A$, and for each $p\in T(x)$ some power $p^{k_p}x$ lies in $N_{(p)}$; so $s=\prod_{p\in T(x)}p^{k_p}\in S_A$ gives $sx\in M$. Hence $M_A=N_A$.
\end{proof}

\begin{proposition}\label{p:obstr}
In ZFC, let $(T_p)_{p\in\Pp}$ be abelian groups, $T_p$ being $p$-primary, for instance $T_p=\Z/p$, and $T=\bigoplus_pT_p$, so that $T_A=\bigoplus_{p\notin A}T_p$. Call a family $g=(g_A)_{A\in U}$, with $g_A\in\prod_{p\notin A}T_p$, \emph{coherent} if $g_A-g_B$ is finitely supported on $A^c\cap B^c$ for all $A,B$, and \emph{trivial} if there is $k\in\prod_pT_p$ with $g_A=k$ on $A^c$ up to finitely many primes, for every $A$. For coherent $g$, put $t_{AB}=(g_A-g_B)|_{A^c\cap B^c}\in T_{A\cup B}$, a $1$-cocycle, and
\[
N_A=\Z_A\oplus T_A,\qquad\varphi_{AB}(r,t)=(r,\,t+r\,t_{AB}).
\]
This is a descent datum, an extension of $f^*\mathcal O_Y$ by $f^*T$, and it is effective if and only if $g$ is trivial.
\end{proposition}

\begin{proof}
The cocycle condition is $t_{AC}=t_{AB}+t_{BC}$ on $A^c\cap B^c\cap C^c$, which holds by definition of $t$. Suppose that $(M,\alpha_A)$ descends the datum, with $\alpha_A\colon M_A\cong N_A$. The projections $N_A\to\Z_A$ are compatible with the $\varphi$'s, so by the equaliser condition \cite[0EUF]{stacks} for $\Z$ they define $\pi\colon M\to\Z$. Then $\operatorname{coker}(\pi)_A=0$ for all $A$, so $\operatorname{coker}\pi=0$, by the equaliser condition again. Choose $m\in M$ with $\pi(m)=1$. Then $\alpha_A(m)=(1,h_A)$ with $h_A\in T_A$, hence finitely supported. Compatibility with the $\varphi$'s gives $g_A+h_A=g_B+h_B$ on $A^c\cap B^c$. These glue to some $k\in\prod_pT_p$, since every $p$ lies in some $A^c$, for instance for $A=\{p\}^c$; and $g_A=k-h_A$ agrees with $k$ on $A^c$ up to finitely many primes.

Conversely, if $g$ is trivial through $k$, put $h_A=(k-g_A)|_{A^c}\in T_A$, $M=\Z\oplus T$ and $\alpha_A(r,t)=(r,\,t+r\,h_A)$. Then $\varphi_{AB}\alpha_A$ and $\alpha_B$ agree over $\Z_{A\cup B}$, because $h_A+t_{AB}=h_B$ on $A^c\cap B^c$.
\end{proof}

\begin{proposition}\label{p:dev}
In ZFC, $f$ is an effective descent morphism for quasi-coherent sheaves if and only if, for every family $(T_p)$ of $p$-primary groups, every coherent family in the sense of Proposition~\ref{p:obstr} is trivial.
\end{proposition}

\begin{proof}
\emph{Only if} is Proposition~\ref{p:obstr}. Conversely, $\Des(f)$ is abelian, with kernels, cokernels and colimits computed componentwise, localisation being exact and commuting with colimits; $f^*$ is exact, preserves colimits and is fully faithful (Theorem~\ref{t:A}). So the effective data, which form the essential image of $f^*$, are closed under cokernels of maps between them, finite direct sums and filtered colimits.

Let $(N_A,\varphi)$ be a datum. Torsion submodules commute with localisation, so the $(N_A)_{\mathrm{tors}}$ form a torsion datum and the quotients a torsion-free one. By Proposition~\ref{p:tt} they are $f^*M'$ and $f^*M''$, with $M'=\bigoplus_pT_p$, $T_p$ being $p$-primary, and $M''$ torsion-free. Write $M''$ as the directed union of its finitely generated submodules $F$, which are free. The datum is the filtered colimit of its pull-backs to the $f^*F$, so we may assume $M''=\Z^n$. For an extension $0\to f^*M'\to N\to f^*\Z^n\to0$, let $E_m$ be its pull-back along the $m$-th inclusion $f^*\Z\to f^*\Z^n$. The sum map $\bigoplus_mE_m\to N$ is onto, and its kernel consists of the $(t_m)\in(f^*M')^n$ with $\sum_mt_m=0$, which is $f^*$ of a submodule of $M'^n$. So $N$ is effective once the $E_m$ are.

So let $0\to f^*M'\to E\to f^*\Z\to0$. Choose $u_A\in E_A$ over $1\in\Z_A$. Then $\varphi_{AB}(u_A)=u_B+t_{AB}$ with $t_{AB}\in M'_{A\cup B}$, and $t_{AC}=t_{AB}+t_{BC}$. Put $B_p=\{p\}^c\in U$, and for $p\notin A$ let $g_A(p)$ be $t_{AB_p}\in M'_{A\cup B_p}=T_p$. For $p\notin A\cup B$, the cocycle relation in $M'_{A\cup B\cup B_p}=T_p$ gives $g_A(p)-g_B(p)=t_{AB}(p)$. So $g$ is coherent, $t_{AB}=(g_A-g_B)|_{A^c\cap B^c}$, and $E$ is the datum of Proposition~\ref{p:obstr} built from $g$. As $g$ is trivial, $E$ is effective.
\end{proof}

\subsection{Topological translation}
Let $\omega^*=\beta\Pp\setminus\Pp$, the space of nonprincipal ultrafilters on $\Pp$, whose clopen sets are the $C^*=\{V:C\in V\}$ for $C\subseteq\Pp$, with $C^*=D^*$ exactly when $C$ and $D$ differ by a finite set. View $U$ as a point of $\omega^*$, and let $O=\omega^*\setminus\{U\}$. A set $C^*$ lies in $O$ if and only if $C\notin U$, and $O$ is the union of the compact open sets $(A^c)^*$, $A\in U$.

\begin{lemma}\label{l:clopen}
Coherent families $(g_A)_{A\in U}$ with $g_A\colon A^c\to\{0,1\}$, taken up to changing each $g_A$ on a finite set, correspond to clopen subsets $S$ of $O$, through $S\cap(A^c)^*=g_A^{-1}(1)^*$. The family is trivial if and only if $S$ is the trace on $O$ of a clopen subset of $\omega^*$. A clopen $S$ of $O$ is not such a trace if and only if $U$ lies in the closure of both $S$ and $O\setminus S$.
\end{lemma}

\begin{proof}
A clopen $S\subseteq O$ meets each compact $(A^c)^*$ in a clopen of it, which is $G_A^*$ for some $G_A\subseteq A^c$ determined up to a finite set; the traces agree on $(A^c)^*\cap(B^c)^*=(A^c\cap B^c)^*$, which says that the indicator functions of the $G_A$ form a coherent family. Conversely a coherent family gives the open set $S=\bigcup_AG_A^*$, whose complement $\bigcup_A(A^c\setminus G_A)^*$ in $O$ is open as well. The clopen sets of $\omega^*$ are the $W^*$, and $S=W^*\cap O$ means $G_A=W\cap A^c$ up to finite sets for all $A$, that is, triviality through the indicator function of $W$. Finally, if $U\notin\overline S$ then $S$ is closed in $\omega^*$, hence clopen; if $U\notin\overline{O\setminus S}$ then $S\cup\{U\}$ is clopen. If $U$ lies in both closures, no clopen $W^*$ has trace $S$: if $U\in W^*$, then $W^*$ is a neighbourhood of $U$, so it meets $O\setminus S$; if $U\notin W^*$, then $\omega^*\setminus W^*$ is a neighbourhood of $U$ disjoint from $S$.
\end{proof}

A point $U$ such that $O$ has a clopen $S$ with $U\in\overline S\cap\overline{O\setminus S}$ is a \emph{strong butterfly point} of $\omega^*$, with wings $S$ and $O\setminus S$, in the sense of Fine and Gillman \cite{fg}. Then $O$ is not $C^*$-embedded in $\omega^*$, that is, $\beta O\neq\omega^*$. A $\{0,1\}$-valued family that is not trivial as a family of sets stays non-trivial in $\prod_p\Z/p$: a trivialising $k$ would give the trivialising set $W=\{p:k(p)=1\}$.

\subsection{Butterfly points of character \texorpdfstring{$\aleph_1$}{aleph1}}
The butterfly lemma of Fine and Gillman \cite{fg} says that under the continuum hypothesis every point of $\omega^*$ is a strong butterfly point; Pitz and Suabedissen give its proof \cite[Lemma~3.1]{pitz}. That proof uses the continuum hypothesis only to obtain a clopen neighbourhood base $\{U_\alpha:\alpha<\omega_1\}$ of the point, together with two properties of $\omega^*$ that hold in ZFC: it is an F-space, and nonempty $G_\delta$ sets have nonempty interior. One builds by recursion on $\alpha<\omega_1$ clopen sets $A_\alpha$, $B_\alpha$ not containing the point, with all $A_\alpha\cap B_\beta$ empty, covering $\omega^*\setminus U_\alpha$ and both meeting $U_\alpha$; the wings are $\bigcup A_\alpha$ and $\bigcup B_\alpha$. So, in ZFC, every point of character $\aleph_1$ is a strong butterfly point. Van Douwen's problem~P20 in the problem list of \emph{Topology Proceedings} \cite{pearl} states this as known: a point with a local base of cardinality $\omega_1$ has $\beta(\omega^*\setminus\{p\})\neq\omega^*$. An ultrafilter on $\Pp$ generated by $\aleph_1$ sets is a point of character $\aleph_1$ of $\omega^*$.

\begin{theorem}\label{t:C}
Suppose that some nonprincipal ultrafilter $U$ on $\Pp$ is generated by $\aleph_1$ sets, as is every one under the continuum hypothesis. Then $f\colon\coprod_{A\in U}\Spec\Z_A\to\Spec\Z$ is a universal effective epimorphism of schemes \cite[0EUG]{stacks} and a descent morphism for quasi-coherent sheaves (Theorem~\ref{t:A}), but not an effective descent morphism. The datum of Proposition~\ref{p:obstr} with $T_p=\Z/p$, built from a clopen subset of $\omega^*\setminus\{U\}$ that is not the trace of a clopen of $\omega^*$, is not effective. Consequently $\QCoh$ is a separated prestack, but not a stack, for the canonical topology on schemes, and for the canonical topology on affine schemes.
\end{theorem}

\begin{proof}
The point $U$ has character $\aleph_1$, so it is a strong butterfly point. Lemma~\ref{l:clopen} gives a coherent $\{0,1\}$-valued family that is not trivial, hence not trivial in $\prod_p\Z/p$, and Proposition~\ref{p:obstr} applies. For affine schemes, $\mathcal W_U$ is a covering family of the canonical topology of affine schemes by \cite[0EUA]{stacks} and the equaliser condition \cite[0EUF]{stacks}, and its descent data are those of $f$.
\end{proof}

\section{The question in ZFC}\label{sec:open}

Van Douwen, Kunen and van Mill give a model of Martin's axiom with $\mathfrak c=\aleph_2$ in which $\beta(\omega^*\setminus\{p\})=\omega^*$ for every point $p$ of $\omega^*$ \cite{dkm}. In that model every $\{0,1\}$-valued coherent family is trivial, by Lemma~\ref{l:clopen}: the indicator function of a clopen of $\omega^*\setminus\{U\}$ whose two parts both have $U$ in their closure would not extend continuously to $\omega^*$. So the obstruction used in Theorem~\ref{t:C} vanishes for every $U$. Whether $\mathcal W_U$ is then an effective descent morphism we do not know: by Proposition~\ref{p:dev} one needs triviality for all coefficient systems $(T_p)$ of $p$-primary groups, which the $C^*$-embedding of $\omega^*\setminus\{U\}$ does not obviously give. According to the solution of problem P20 in \cite{pearl}, the equality $\beta(\omega^*\setminus\{p\})=\omega^*$ for every $p$ also holds, by Malykhin, in any model where at least as many Cohen reals are added as there are reals in the ground model; the same question arises there. Also open are whether some $U$ gives effectivity in ZFC, and whether there is a counterexample in ZFC, from this family or from another. Set theory seems to enter the question in an essential way.

\subsection*{What is not done}
Theorem~\ref{t:A} is formalised only in its affine core (Section~\ref{sec:lean}); Theorems~\ref{t:B} and~\ref{t:C} and the propositions of Section~\ref{sec:nqc} are not formalised. None of the arguments has been checked by a specialist of descent or of set-theoretic topology. The status of effectivity along $\mathcal W_U$ in ZFC, and in the model of van Douwen, Kunen and van Mill, is not settled. The quotation from the lecture rests on a transcription that has not been checked by ear.

\subsection*{Use of language models}
The question was located in the recording of the Buffalo lecture, and the literature searched, with language models (Claude Opus~5.5, Anthropic), which also found the arguments of Sections~\ref{sec:A}--\ref{sec:nqc}, wrote the Lean file and a first draft of this text, under the author's direction. The references were opened on their pages; the proofs of Pitz and Suabedissen and of the Stacks Project that are used were read. The Lean file is checked by the Lean kernel, and anyone can rerun it.

\begin{thebibliography}{99}
\bibitem{ab} K.~Aoki, R.~Burklund, Inaccessibility of the flat topology, 2026, \href{https://arxiv.org/abs/2607.09440}{arXiv:2607.09440}.
\bibitem{af} Y.~André, L.~Fiorot, On the canonical, fpqc, and finite topologies on affine schemes. The state of the art, \emph{Ann.\ Sc.\ Norm.\ Super.\ Pisa Cl.\ Sci.~(5)} 23 (2022), 81--114, \href{https://doi.org/10.2422/2036-2145.201912_006}{doi:10.2422/2036-2145.201912\_006}; \href{https://arxiv.org/abs/1912.04957}{arXiv:1912.04957}.
\bibitem{buffalo} A.~Grothendieck, course on algebraic geometry, State University of New York at Buffalo, session of 11 July 1973, recorded on the second video of the course on group schemes, at 6:26:20--6:27:24 of \url{https://www.youtube.com/watch?v=3ULFIl_2Wis&t=6h26m20s}.
\bibitem{dkm} E.~K.~van Douwen, K.~Kunen, J.~van Mill, There can be $C^*$-embedded dense proper subspaces in $\beta\omega-\omega$, \emph{Proc.\ Amer.\ Math.\ Soc.} 105 (1989), 462--470, \href{https://doi.org/10.1090/S0002-9939-1989-0977925-7}{doi:10.1090/S0002-9939-1989-0977925-7}.
\bibitem{fg} N.~J.~Fine, L.~Gillman, Extension of continuous functions in $\beta\mathbf N$, \emph{Bull.\ Amer.\ Math.\ Soc.} 66 (1960), 376--381, \href{https://doi.org/10.1090/S0002-9904-1960-10460-0}{doi:10.1090/S0002-9904-1960-10460-0}.
\bibitem{kern} S.~Kern, The crisp topology, a refinement of the fpqc topology, 2026, \href{https://arxiv.org/abs/2603.25272}{arXiv:2603.25272}.
\bibitem{mesablishvili} B.~Mesablishvili, Pure morphisms of commutative rings are effective descent morphisms for modules---a new proof, \emph{Theory Appl.\ Categ.} 7 (2000), no.~3, 38--42, \url{http://www.tac.mta.ca/tac/volumes/7/n3/7-03abs.html}.
\bibitem{mesablishvili-acs} B.~Mesablishvili, Descent theory for schemes, \emph{Appl.\ Categ.\ Structures} 12 (2004), 485--512, \href{https://doi.org/10.1023/B:APCS.0000049314.33172.0d}{doi:10.1023/B:APCS.0000049314.33172.0d}.
\bibitem{mesablishvili2004} B.~Mesablishvili, More on descent theory for schemes, \emph{Georgian Math.\ J.} 11 (2004), 783--800, \href{https://doi.org/10.1515/GMJ.2004.783}{doi:10.1515/GMJ.2004.783}.
\bibitem{mo127362} O.~Benoist, answer to \emph{Counter-example to faithfully flat descent}, MathOverflow question 127362, 12 April 2013, \url{https://mathoverflow.net/questions/127362}.
\bibitem{olivier} J.-P.~Olivier, Descente par morphismes purs, \emph{C.\ R.\ Acad.\ Sci.\ Paris Sér.~A} 271 (1970), 821--823.
\bibitem{orlov} D.~O.~Orlov, Quasicoherent sheaves in commutative and noncommutative geometry, \emph{Izv.\ Math.} 67 (2003), 535--554, \href{https://doi.org/10.1070/IM2003v067n03ABEH000437}{doi:10.1070/IM2003v067n03ABEH000437}.
\bibitem{pearl} E.~Pearl (ed.), Problems from \emph{Topology Proceedings}, Topology Atlas, 2003, \href{https://arxiv.org/abs/math/0312456}{arXiv:math/0312456}.
\bibitem{picavet} G.~Picavet, Submersion et descente, \emph{J.\ Algebra} 103 (1986), 527--591, \href{https://doi.org/10.1016/0021-8693(86)90152-3}{doi:10.1016/0021-8693(86)90152-3}.
\bibitem{pitz} M.~F.~Pitz, R.~Suabedissen, Compactifications of $\omega^*\setminus\{x\}$ and $S_\kappa\setminus\{x\}$, \emph{Topology Proc.} 45 (2015), 151--173, zbMATH 1315.54023; \href{https://arxiv.org/abs/1310.0678}{arXiv:1310.0678} (numbering of the arXiv version).
\bibitem{repo} Commutator-IO, grothendieck-project, \url{https://github.com/Commutator-IO/grothendieck-project}: \texttt{lean/Buffalo/QcohDescent.lean}, and the index of results of the Buffalo courses (\texttt{docs/buffalo-ag/index.tex}).
\bibitem{rydh} D.~Rydh, Submersions and effective descent of étale morphisms, \emph{Bull.\ Soc.\ Math.\ France} 138 (2010), 181--230, \href{https://doi.org/10.24033/bsmf.2588}{doi:10.24033/bsmf.2588}.
\bibitem{stacks} The Stacks Project Authors, \emph{The Stacks Project}, \url{https://stacks.math.columbia.edu}, tags as cited.
\end{thebibliography}
\end{document}
